\section{问题1：完全信息下的日计划购电模型}

\subsection{模型建立}

问题1 中电价、小区负载与光伏预测功率全天已知，故可在 0:00 一次求解全天最优计划。以 \cref{eq:balance}--\cref{eq:terminal} 为基础，建立如下线性规划：

\begin{equation}
  \min_{\{g_t,a_t,b_t,s_t,E_t\}}\ F_1=\sum_{t=1}^{T}p_tg_t\text{。}
  \label{eq:p1-obj}
\end{equation}

\begin{equation}
  \text{s.t.}\quad
  \begin{cases}
    g_t+\eta_db_t+P_t\Delta t=L_t\Delta t+\dfrac{a_t}{\eta_c}+s_t, & \text{母线功率平衡}\\[4pt]
    E_t=E_{t-1}+a_t-b_t, & \text{储能动态}\\[4pt]
    E_{144}=E_0=6000, & \text{日周期条件}\\[4pt]
    1200\le E_t\le10800,\quad 0\le a_t,b_t\le 833.33, & \text{状态与功率约束}\\[4pt]
    g_t\ge0,\quad s_t\ge0 . & \text{非负与弃光}
  \end{cases}\text{。}
  \label{eq:p1-cons}
\end{equation}

模型为线性规划：决策变量 $720$ 个（$g,a,b,s$ 各 $144$ 个、$E$ 为 $144$ 个），等式约束 $289$ 条（平衡 $144$ 条、动态 $144$ 条、终端 $1$ 条），状态与功率约束以变量界表达。采用 HiGHS 求解器\cite{huangfu2018highs} 与 SciPy 科学计算库\cite{virtanen2020} 求解，规模极小，可在毫秒级得到全局最优解。

\subsection{模型性质与最优计划结构}

\textbf{（1）全天购电量分解。} 由 \cref{eq:identity}，问题1 的全天购电量满足"全天购电量 $=$ 净缺口 $+$ 弃光 $+0.2111A$"。代入实测数据：净缺口 $55541.97$ \si{kWh}，弃光 $0$，电池侧吞吐量 $A=18666.60$ \si{kWh}，理论购电量为 $55541.97+\frac{19}{90}\times18666.60=59482.70$ \si{kWh}，与求解结果一致（残差 $2.2\times10^{-11}$）；其中 $A=18666.60$ \si{kWh} 为电池侧吞吐量，等于 $0.9\times$ 母线侧充电量（$0.9\times20740.67=18666.60$ \si{kWh}）。该恒等式说明储能套利的"代价"是 $23.46\%$ 的能量损耗，其收益必须来自足够大的时段价差。

\textbf{（2）套利判据。} 设电池侧电量 $A$ 在时段 $t_c$ 充入、在时段 $t_d$ 放出，则费用变化为

\begin{equation}
  \Delta F=A\left(\frac{p_{t_c}}{\eta_c}-\eta_d\,p_{t_d}\right),
  \label{eq:arb}
\end{equation}

当且仅当 $p_{t_d}/p_{t_c}>1/(\eta_c\eta_d)=1.235$ 时值得搬运。用附件1 电价核对主要窗口：夜间 $0.42$ 元/kWh 充电、晚峰 $1.30$ 元/kWh 放电的价比为 $3.10$；午间 11 时 $0.47$ 充电、晚峰 $1.30$ 放电的价比为 $2.77$；清晨 4 时 $0.44$、早峰 8 时 $1.08$ 的价比为 $2.46$。三者均远超 $1.235$，故最优解必然表现为"夜充、峰放、日末回充"的结构。

\textbf{（3）阈值策略验证。} \cref{eq:threshold} 预测了"充电时段电价低于阈值、放电时段电价高于阈值"的分布特征。\cref{fig:p1-threshold} 将 $144$ 个时段的充放电状态标注在电价轴上：充电时段电价集中在 $0.37$--$0.53$ 元/kWh（夜间与午间谷段），放电时段电价集中在 $0.87$--$1.35$ 元/kWh（早晚峰），两个区间几乎无重叠，与理论阈值策略一致。这说明最优计划不是"逐时段贪心"，而是由储电量影子价格 $\lambda$ 统一决定的套利行为。

\begin{figure}[H]
  \centering
  \includegraphics[width=0.66\linewidth]{fig04.png}
  \caption{问题1 的阈值策略验证：(a) 附件1 电价曲线与阈值示意；(b) 充、放、闲置状态在电价轴上的分布}
  \label{fig:p1-threshold}
\end{figure}

\subsection{求解结果}

求解得全天购电费 $F_1^\star=35101.57$ 元、全天购电量 $59482.70$ \si{kWh}；若不计储能（逐时段按净缺口购电），全天购电费为 $48052.05$ 元，故储能带来 $27.0\%$ 的费用下降，即 $12950.48$ 元。储电量轨迹为：0:00 与 24:00 均为 $6000$ \si{kWh}，最低 $1200$ \si{kWh}（20:40），最高 $10800$ \si{kWh}（05:40），全程满足 $[1200,10800]$ \si{kWh} 的约束；电池侧单时段最大充电量达到上限 $833.33$ \si{kWh}（功率约束起作用），最大放电量为 $795.17$ \si{kWh}（电池侧；对应母线侧 $0.9\times795.17=715.65$ \si{kWh}）。

\begin{table}[H]
  \centering
  \caption{问题1 指定时段的购电量与全天合计（表1 格式）}
  \label{tab:p1-t1}
  \small
  \footnotesize\begin{tabular}{cc|cc|cc}
    \toprule[1.5pt]
    时间段 & 购电量 (\si{kWh}) & 时间段 & 购电量 (\si{kWh}) & 时间段 & 购电量 (\si{kWh}) \\
    \midrule[1pt]
    10:00--10:10 & 0.00 & 12:00--12:10 & 573.01 & 14:00--14:10 & 0.00 \\
    16:00--16:10 & 445.43 & 18:00--18:10 & 531.89 & 20:00--20:10 & 0.00 \\
    \midrule[1pt]
    \multicolumn{5}{l}{全天购电量 / 全天购电费} & 59482.70 \si{kWh} / 35101.57 元 \\
    \bottomrule[1.5pt]
  \end{tabular}
\end{table}

\begin{table}[H]
  \centering
  \caption{问题1 储能设备在指定时间段的充放电量与储电量（表2 格式）}
  \label{tab:p1-t2}
  \small
  \begin{tabular}{lcc}
    \toprule[1.5pt]
    时间段 & 充电量 (\si{kWh}) & 放电量 (\si{kWh}) \\
    \midrule[1pt]
    0:00--4:00   & 4407.41 & 0.00 \\
    4:00--8:00   & 925.93  & 6365.84 \\
    8:00--12:00  & 4417.59 & 1703.00 \\
    12:00--16:00 & 5656.41 & 91.10 \\
    16:00--20:00 & 0.00    & 5780.13 \\
    20:00--24:00 & 5333.33 & 2859.87 \\
    \midrule[1pt]
    0:00 / 24:00 储电量 & \multicolumn{2}{c}{6000.00 \si{kWh} / 6000.00 \si{kWh}} \\
    \bottomrule[1.5pt]
  \end{tabular}
\end{table}

\cref{tab:p1-t1} 给出表1 指定时段的购电量，\cref{tab:p1-t2} 给出储能充放电量。由 \cref{tab:p1-t2,fig:p1-plan} 可见，储能全天完成约 1.9 个充放循环：0:00--4:00 夜间低谷充电 $4407.41$ \si{kWh}，4:00--8:00 早峰放电 $6365.84$ \si{kWh}；8:00--12:00 与 12:00--16:00 在午间次低谷继续充电（合计 $10074.00$ \si{kWh}），16:00--20:00 晚峰放电 $5780.13$ \si{kWh}；20:00--24:00 电价回落，再充入 $5333.33$ \si{kWh} 以满足日周期条件。充电量按母线侧口径统计，故全天充电量（$20740.67$ \si{kWh}）大于放电量（$16799.94$ \si{kWh}），差额即往返效率损耗。

\begin{figure}[H]
  \centering
  \includegraphics[width=0.66\linewidth]{fig05.png}
  \caption{问题1 的最优计划：(a) 计划购电功率与净缺口；(b) 储电量轨迹（日周期 $E_{144}=E_0$）}
  \label{fig:p1-plan}
\end{figure}

\subsection{灵敏度分析}

为检验结论的稳健性，用参数化线性规划对效率、功率上限、储电量上下限、初始电量与电价水平作单因素扰动，结果见 \cref{fig:p1-sens} 与 \cref{tab:p1-sens}。

\begin{figure}[H]
  \centering
  \includegraphics[width=0.66\linewidth]{fig06.png}
  \caption{问题1 的单因素灵敏度：(a) 各扰动情景下的全天购电费；(b) 龙卷风图（按影响大小排序）}
  \label{fig:p1-sens}
\end{figure}

\begin{table}[H]
  \centering
  \caption{问题1 灵敏度分析结果}
  \label{tab:p1-sens}
  \small
  \begin{tabular}{lccc}
    \toprule[1.5pt]
    扰动情景 & 全天购电费 (元) & 相对基准变化 & 电池侧吞吐量 (\si{kWh}) \\
    \midrule[1pt]
    基准                      & 35101.57 & ---       & 18666.60 \\
    效率 $\eta=0.85$          & 36548.22 & $+4.12\%$ & 19199.93 \\
    效率 $\eta=0.95$          & 33756.16 & $-3.83\%$ & 18811.66 \\
    最大功率 $4000$ \si{kW}   & 35327.90 & $+0.64\%$ & 18510.81 \\
    最大功率 $6000$ \si{kW}   & 35056.72 & $-0.13\%$ & 18769.48 \\
    储量下限 $2400$ \si{kWh}  & 36014.15 & $+2.60\%$ & 16901.22 \\
    储量上限 $9600$ \si{kWh}  & 36008.47 & $+2.58\%$ & 16901.22 \\
    初始电量 $3000$ \si{kWh}  & 35125.20 & $+0.07\%$ & 18666.60 \\
    初始电量 $9000$ \si{kWh}  & 35093.43 & $-0.02\%$ & 18666.60 \\
    电价 $+10\%$              & 38611.72 & $+10.00\%$ & 18666.60 \\
    电价 $-10\%$              & 31591.41 & $-10.00\%$ & 18666.60 \\
    \bottomrule[1.5pt]
  \end{tabular}
\end{table}

三点结论：其一，\textbf{效率是最敏感的内部参数}——往返效率由 $81\%$ 降至 $72.25\%$（$\eta=0.85$）时费用上升 $4.12\%$，且吞吐量反而增加（需搬运更多电量补偿损耗），说明效率直接决定套利的可行空间；其二，\textbf{储电量上下限的影响（约 $2.6\%$）大于功率上限（约 $0.6\%$）}，因为可用容量仅为日净缺口的 $17\%$，容量而非功率是套利的紧约束；其三，\textbf{日周期条件下初始电量几乎无影响}（$\pm0.07\%$），因为 $E_{144}=E_0$ 要求全天自平衡，初始电量只改变循环起点而不改变搬运总量；电价整体升降则近似线性传递到费用（弹性约为 1）。
